\documentclass[11pt]{article}

% ============================================================
%  The SRDS Formal Framework
%  A self-contained mathematical theory of critical transitions
%  in self-referential dissipative systems.
%
%  Standalone document. No external repository references.
%  Only established scholarly literature is cited.
% ============================================================

\usepackage[utf8]{inputenc}
\usepackage{lmodern}
\usepackage[T1]{fontenc}
\usepackage{amsmath,amssymb,amsthm,mathtools}
\usepackage[margin=1in]{geometry}
\IfFileExists{microtype.sty}{\usepackage[protrusion=true,expansion=false,final]{microtype}}{}
\tolerance=2000
\emergencystretch=3em
\usepackage{enumitem}
\usepackage[numbers]{natbib}
\usepackage{hyperref}
\setcounter{tocdepth}{1}
\hypersetup{colorlinks=true,linkcolor=black,citecolor=black,urlcolor=black}

\theoremstyle{plain}
\newtheorem{theorem}{Theorem}[section]
\newtheorem{lemma}[theorem]{Lemma}
\newtheorem{proposition}[theorem]{Proposition}
\newtheorem{corollary}[theorem]{Corollary}
\theoremstyle{definition}
\newtheorem{definition}[theorem]{Definition}
\newtheorem{axiom}{Axiom}
\renewcommand{\theaxiom}{A\arabic{axiom}}
\setcounter{axiom}{-1}
\newtheorem{regularity}{Regularity Condition}
\renewcommand{\theregularity}{L\arabic{regularity}}
\newtheorem{assumption}{Assumption}
\renewcommand{\theassumption}{R\arabic{assumption}}
\newtheorem{conjecture}[theorem]{Conjecture}
\newtheorem{hypothesis}[theorem]{Hypothesis}
\theoremstyle{remark}
\newtheorem{remark}[theorem]{Remark}

\DeclareMathOperator{\Var}{Var}
\DeclareMathOperator{\Cov}{Cov}
\DeclareMathOperator{\Tr}{Tr}
\DeclareMathOperator{\sgn}{sign}
\DeclareMathOperator*{\argmax}{arg\,max}
\DeclareMathOperator*{\argmin}{arg\,min}
\newcommand{\R}{\mathbb{R}}
\newcommand{\N}{\mathbb{N}}
\newcommand{\E}{\mathbb{E}}
\newcommand{\norm}[1]{\left\lVert#1\right\rVert}
\newcommand{\abs}[1]{\left\lvert#1\right\rvert}
\newcommand{\ip}[2]{\langle #1, #2\rangle}
\newcommand{\Rc}{\mathcal R}
\newcommand{\Ac}{\mathcal A}

\title{\textbf{The SRDS Formal Framework}\\[0.3em]
\large A Self-Contained Mathematical Theory of Critical Transitions\\ in Self-Referential Dissipative Systems}
\author{Ferdinand Maria Scheßl}
\date{2026}

\begin{document}
\maketitle

\begin{abstract}
We develop, in full and self-contained form, the mathematical theory of \emph{self-referential dissipative systems} (SRDS): sequences generated by a state recursion whose update operator is itself driven by the state it produces. The class is fixed by five axioms (self-reference, dissipation, accumulation, transition, irreversibility) together with three regularity conditions, and we prove that accumulation, transition, and irreversibility are \emph{theorems} reducible to self-reference and dissipation under the regularity conditions. On this base we prove three structural results: an information--stability correspondence relating the sign of the maximal Lyapunov exponent to the trend of mutual information (T1); an autopoietic stationarity theorem giving a unique, globally attracting non-equilibrium steady state of the operator recursion (T2); and a Goldilocks principle locating optimal function in a neighbourhood of, rather than exactly at, criticality (T3). We construct the two-scale renormalization operator that maps the dynamics across observational scales, prove the existence and uniqueness of its monotone fixed point through Birkhoff--Hopf projective contraction, and identify the cumulative inflection as the order parameter of the class. Finally we prove a No-Go theorem: within reparametrization-closed operator classes no universal numerical inflection location exists; the symmetric value is recovered exactly when the kernel is mirror-symmetric, and is otherwise a substrate-specific measured quantity. We close with an elimination ledger making precise what cannot be concluded without additional structure, and with the spectral and contraction substructure governing convergence rates. Every result is proved in full, reduced to a cited result of the established literature, or stated under explicitly named additional hypotheses, with the conditions on which it rests made explicit. The framework is metric-agnostic and is instantiated by, but does not depend on, any particular measurement system.
\end{abstract}

\tableofcontents

% ============================================================
\section{Preliminaries and Notation}\label{sec:prelim}
% ============================================================

This section fixes the analytic and information-theoretic vocabulary used throughout and records the standing regularity assumptions. All statements are proved or reduced to standard references.

\subsection{Function spaces}

Let $C([0,1])$ denote the space of continuous real functions on $[0,1]$ with the supremum norm $\norm{Q}_\infty := \sup_{x\in[0,1]}\abs{Q(x)}$.

\begin{proposition}[Completeness]\label{prop:complete}
$(C([0,1]),\norm{\cdot}_\infty)$ is a complete metric space.
\end{proposition}
\begin{proof}
Let $(Q_n)_{n\in\N}$ be Cauchy in $\norm{\cdot}_\infty$. For each fixed $x$, $\abs{Q_n(x)-Q_m(x)}\le\norm{Q_n-Q_m}_\infty$, so $(Q_n(x))_n$ is Cauchy in $\R$ and, by completeness of $\R$, converges to a limit we call $Q(x)$. The convergence is uniform: given $\varepsilon>0$ choose $N$ with $\norm{Q_n-Q_m}_\infty<\varepsilon$ for $n,m\ge N$; letting $m\to\infty$ in $\abs{Q_n(x)-Q_m(x)}<\varepsilon$ gives $\abs{Q_n(x)-Q(x)}\le\varepsilon$ for all $x$ and all $n\ge N$, i.e.\ $\norm{Q_n-Q}_\infty\le\varepsilon$. A uniform limit of continuous functions is continuous: for $n\ge N$ and $\abs{x-x_0}$ small enough that $\abs{Q_n(x)-Q_n(x_0)}<\varepsilon$ (continuity of $Q_n$), the triangle inequality $\abs{Q(x)-Q(x_0)}\le\abs{Q(x)-Q_n(x)}+\abs{Q_n(x)-Q_n(x_0)}+\abs{Q_n(x_0)-Q(x_0)}<3\varepsilon$ holds. Hence $Q\in C([0,1])$ and $Q_n\to Q$.
\end{proof}

Let $BV([0,1])$ denote the functions of bounded variation, with total variation
\[
\Var(Q):=\sup_{P}\sum_{i}\abs{Q(x_i)-Q(x_{i-1})}
\]
over partitions $P=\{0=x_0<\dots<x_k=1\}$.

\begin{proposition}\label{prop:bv}
Every $Q\in BV([0,1])$ is bounded, with $\abs{Q(x)}\le\abs{Q(0)}+\Var(Q)$ for all $x$.
\end{proposition}
\begin{proof}
For fixed $x$ the partition $\{0,x\}$ gives $\abs{Q(x)-Q(0)}\le\Var(Q)$, whence $\abs{Q(x)}\le\abs{Q(0)}+\abs{Q(x)-Q(0)}\le\abs{Q(0)}+\Var(Q)$.
\end{proof}

We write $\mathcal{M}\subset BV([0,1])$ for the set of \emph{anchored monotone profiles}: nondecreasing $Q$ with $Q(0)=0$, $Q(1)=1$. Such a $Q$ is the cumulative distribution of a probability measure on $[0,1]$; by Proposition~\ref{prop:bv}, $\mathcal{M}$ is uniformly bounded.

\subsection{Information theory}

\begin{definition}
For a discrete random variable $X$ with law $p$, the Shannon entropy is $H(X):=-\sum_x p(x)\log p(x)$ (convention $0\log 0=0$). For a pair $(X,Y)$ the joint entropy is $H(X,Y):=-\sum_{x,y}p(x,y)\log p(x,y)$, the conditional entropy $H(X\mid Y):=H(X,Y)-H(Y)$, and the mutual information $I(X;Y):=H(X)+H(Y)-H(X,Y)$.
\end{definition}

\begin{lemma}[Properties of mutual information]\label{lem:mi}
$I$ satisfies: (i) $I(X;Y)\ge 0$; (ii) $I(X;Y)=I(Y;X)$; (iii) $I(X;Y)=0$ iff $X,Y$ are independent; (iv) $I(X;Y)\le\min\{H(X),H(Y)\}$.
\end{lemma}
\begin{proof}
Write $I(X;Y)=\sum_{x,y}p(x,y)\log\frac{p(x,y)}{p(x)p(y)}$. (i) By Jensen's inequality applied to the convex function $t\mapsto t\log t$, equivalently by the nonnegativity of the Kullback--Leibler divergence $D\!\left(p(x,y)\,\|\,p(x)p(y)\right)=I(X;Y)\ge 0$ \citep{coverthomas2006}. (ii) Symmetry is immediate from $I=H(X)+H(Y)-H(X,Y)$. (iii) Equality $I=0$ holds iff $D=0$, i.e.\ $p(x,y)=p(x)p(y)$ everywhere, which is independence. (iv) $I(X;Y)=H(X)-H(X\mid Y)$ and $H(X\mid Y)\ge 0$ for discrete variables give $I\le H(X)$; symmetry gives $I\le H(Y)$.
\end{proof}

\begin{theorem}[Data-processing inequality]\label{thm:dpi}
If $X\to Y\to Z$ is a Markov chain then $I(X;Z)\le I(X;Y)$.
\end{theorem}
\begin{proof}
Standard; see \citep[Thm.~2.8.1]{coverthomas2006}. Processing of $Y$ cannot increase information about $X$.
\end{proof}

For continuous laws we use the differential entropy $h(X):=-\int p\log p$ and the Gaussian maximum-entropy bound: among densities on $\R^n$ with covariance $\Sigma$, the Gaussian maximizes $h$, with $h(X)\le\tfrac12\log\!\big((2\pi e)^n\det\Sigma\big)$ \citep[Thm.~8.6.5]{coverthomas2006}.

\subsection{The measurement model}

The framework is metric-agnostic: it constrains the \emph{geometry of a trajectory}, not the instrument that records it. We fix the minimal measurement vocabulary.

\begin{definition}[Specimen and deformation curve]
A \emph{specimen} is a finite ordered sequence $(z_1,\dots,z_T)$ of states $z_t\in\R^{m}$, obtained from an embedding map $\varphi$ of the raw sequence into $\R^m$. The index $t$ is ordered loading time, not an i.i.d.\ sample index; the ordered sequence is the \emph{deformation curve} of the specimen. For window size $W$ and $t\ge W$ the window is $\mathcal{W}(t):=\{t-W+1,\dots,t\}$ (with $\mathcal{W}(t):=\{1,\dots,t\}$ for $t<W$).
\end{definition}

Because the index is loading time, successive measurements along a deformation curve are \emph{autocorrelated by construction}; this is signal, not nuisance, and inference at the specimen level must treat the whole curve as the unit of analysis (Section~\ref{sec:operational}).

\subsection{Standing regularity assumptions}

The dynamical results invoke the following conditions. They are sufficient, not necessary, and each is satisfied by standard generative models of ordered embedded sequences.

\begin{assumption}[Stochastic flow]\label{R:flow}
The state sequence $(z_t)$ is generated by a discrete-time process whose continuous-time interpolation solves $dz=F(z)\,dt+\sigma\,dW_t$, with $F:\R^m\to\R^m$ locally Lipschitz, $\sigma>0$, and $W_t$ standard Brownian motion in $\R^m$.
\end{assumption}

\begin{assumption}[Lyapunov spectrum]\label{R:lyap}
Along trajectories the variational equation $\dot\xi=DF(z(t))\,\xi$ admits a maximal Lyapunov exponent $\lambda_{\max}:=\lim_{t\to\infty}\tfrac1t\log\norm{\xi(t)}/\norm{\xi(0)}$, which exists for a.e.\ initial condition by the multiplicative ergodic theorem \citep{oseledets1968}.
\end{assumption}

\begin{assumption}[Bi-Lipschitz embedding]\label{R:embed}
On the support of the trajectory law, $\varphi$ is bi-Lipschitz: there are $0<c_1\le c_2<\infty$ with $c_1\,d_{\mathcal X}(s,s')\le\norm{\varphi(s)-\varphi(s')}\le c_2\,d_{\mathcal X}(s,s')$ for a base metric $d_{\mathcal X}$.
\end{assumption}

\begin{assumption}[Ergodicity]\label{R:ergodic}
$(z_t)$ is stationary and ergodic with invariant law $\mu$ having log-concave density on $\R^m$; hence time averages converge to $\mu$-averages \citep{birkhoff1931} and concentration holds \citep{bobkov2003}.
\end{assumption}

\begin{assumption}[Finite information]\label{R:mi}
For every $t$ and every finite measurable partition, the discretized state coordinates have finite Shannon entropy, so $I(\cdot;\cdot)$ is well defined.
\end{assumption}

% ============================================================
\section{The SRDS Axiomatic Class}\label{sec:axioms}
% ============================================================

\subsection{Two recursions}

\begin{definition}[Self-referential system]\label{def:srds}
A self-referential dissipative system is a pair of coupled recursions
\begin{equation}\label{eq:rec}
z_{t+1}=F_t(z_t,u_t),\qquad F_{t+1}=\Phi(F_t,z_t,u_t),
\end{equation}
where $z_t\in\R^m$ is the state, $u_t$ an exogenous input, $F_t\in\mathcal F$ the update operator drawn from a complete metric space $(\mathcal F,d_{\mathcal F})$, and $\Phi:\mathcal F\times\R^m\times\mathcal U\to\mathcal F$ the meta-operator. A configuration $F^*$ is \emph{autopoietic} if $\Phi(F^*,z^*,u^*)=F^*$.
\end{definition}

The first recursion is classical feedback dynamics \citep{wiener1948,ashby1956}. It is insufficient alone: a specimen remembers its loading history (the Bauschinger effect), a sequence carries path dependence. The second recursion makes that memory explicit: the operator that writes the next state is itself rewritten during writing. This is the operative closure of \citep{maturana1980}.

\subsection{The five axioms}

We collect the cumulative dissipation along a trajectory. Under Assumption~\ref{R:flow} let $V:\mathcal F\to[0,\infty)$ be a measurable functional and define the one-step expected decrement and the cumulative damage
\begin{equation}\label{eq:damage}
\delta_t:=V(F_t)-\E\!\left[V(F_{t+1})\mid \mathcal G_t\right],\qquad
\Sigma_t:=\sum_{s<t}\delta_s,
\end{equation}
where $(\mathcal G_t)$ is the natural filtration.

\begin{axiom}[Self-reference]\label{ax:self}
The update operator evolves by the second recursion of \eqref{eq:rec}: $F_{t+1}=\Phi(F_t,z_t,u_t)$; the operator is itself a state, rewritten as it acts.
\end{axiom}

\begin{axiom}[Dissipation]\label{ax:diss}
There is a Lyapunov functional $V\ge0$ with $\E[V(F_{t+1})\mid\mathcal G_t]\le V(F_t)$, i.e.\ $\delta_t\ge0$ in conditional expectation. Equivalently the system exports entropy to maintain bounded internal organization far from equilibrium \citep{prigogine1977}.
\end{axiom}

\begin{axiom}[Accumulation]\label{ax:acc}
The damage $\Sigma_t$ is monotone nondecreasing: $\Sigma_{t+1}\ge\Sigma_t$.
\end{axiom}

\begin{axiom}[Transition]\label{ax:trans}
The normalized damage $\sigma(t):=\Sigma_t/\Sigma_\infty$ admits a unique inflection point $\hat a$ (sigmoid form), the order parameter of the trajectory \citep{scheffer2009}.
\end{axiom}

\begin{axiom}[Irreversibility]\label{ax:irr}
A trajectory that has crossed the inflection ($\Sigma>\hat a$) returns to its pre-transition configuration with probability strictly less than one \citep{seifert2012}.
\end{axiom}

\subsection{Regularity conditions}

\begin{regularity}[Lyapunov half-order]\label{L1}
$V$ is bounded below (normalized $V\ge0$) and adapted to $(\mathcal G_t)$, so the increments $\delta_t$ are nonnegative adapted random variables and $\Sigma_t$ is a nondecreasing adapted process.
\end{regularity}

\begin{regularity}[Saturation bound]\label{L2}
The expected dissipation rate $t\mapsto\E[\delta_t]$ is unimodal: it has a single interior maximum, rising from a quasi-elastic onset and decaying as the trajectory approaches the autopoietic configuration. (Finiteness of the cumulative dissipation, $\E[\Sigma_\infty]\le V(F_0)$, follows from $V\ge0$ in \ref{L1}.)
\end{regularity}

\begin{regularity}[Asymmetry measure]\label{L3}
The stochastic-thermodynamic entropy production $\Delta\Sigma$ along any segment satisfies a detailed fluctuation relation $\mathbb P[\gamma^{\dagger}]/\mathbb P[\gamma]=e^{-\Delta\Sigma}$, where $\gamma^{\dagger}$ is the time-reversed segment of $\gamma$ \citep{crooks1999,seifert2012}.
\end{regularity}

\subsection{The reduction theorem}

The point of the axiomatization is that only \ref{ax:self}--\ref{ax:diss} are primitive: accumulation, transition, and irreversibility are \emph{theorems} once the regularity conditions hold. Making the reduction explicit lets the load-bearing capacity of each axiom be tested on each substrate separately, and it is exactly this structure that produces the predicted null findings of a substrate that violates one regularity condition.

\begin{theorem}[Reduction]\label{thm:reduction}
Under Axioms \ref{ax:self}--\ref{ax:diss}:
\begin{enumerate}[label=(\roman*),leftmargin=2.2em]
\item \emph{(}Accumulation\emph{)} \ref{L1}+\ref{L2} imply \ref{ax:acc}, and $\Sigma_t\uparrow\Sigma_\infty<\infty$.
\item \emph{(}Transition\emph{)} \ref{L2} implies \ref{ax:trans}: $\sigma(t)=\Sigma_t/\Sigma_\infty$ is sigmoid with a unique inflection $\hat a=\argmax_t \E[\delta_t]$.
\item \emph{(}Irreversibility\emph{)} \ref{L3} implies \ref{ax:irr}.
\end{enumerate}
\end{theorem}
\begin{proof}
\textbf{(i)} By Axiom~\ref{ax:diss} the conditional increments satisfy $\delta_t\ge0$; by \ref{L1} they are nonnegative and adapted, so $\Sigma_{t+1}=\Sigma_t+\delta_t\ge\Sigma_t$ pathwise, which is Axiom~\ref{ax:acc}; thus $\Sigma_t$ is the predictable compensator of the nonnegative supermartingale $V(F_t)$. Taking expectations and telescoping, $\E[\Sigma_t]=\sum_{s<t}\E[\delta_s]=V(F_0)-\E[V(F_t)]\le V(F_0)<\infty$, using $V\ge0$ (\ref{L1}). A nondecreasing process with uniformly bounded expectation converges almost surely by the monotone convergence theorem to a limit $\Sigma_\infty$ with $\E[\Sigma_\infty]\le V(F_0)<\infty$; hence $\Sigma_\infty<\infty$ a.s.

\textbf{(ii)} Work with the mean damage curve $s(t):=\E[\Sigma_t]/\E[\Sigma_\infty]$, the deterministic profile actually fitted. By (i) it increases from $0$ to $1$, with increments $s(t+1)-s(t)=\E[\delta_t]/\E[\Sigma_\infty]=:r(t)\ge0$. By \ref{L2}, $r$ is unimodal with a single interior maximum at $\hat a:=\argmax_t r(t)$. Hence the second difference $s(t+1)-2s(t)+s(t-1)=r(t)-r(t-1)$ changes sign exactly once, from positive to negative, at $\hat a$: $s$ is convex for $t<\hat a$ and concave for $t>\hat a$. A bounded monotone function with exactly one convex-to-concave change is sigmoid with a unique inflection at $\hat a$, which is Axiom~\ref{ax:trans}. The inflection is the loading time of maximal expected damage rate: the yield point of the trajectory.

\textbf{(iii)} The return event $\mathcal R$, that the time-reversed dynamics restores the pre-transition configuration, is realized exactly by the reversals $\gamma^\dagger$ of the forward crossing paths $\gamma$, those with $\Delta\Sigma(\gamma)>0$ (strict, because a crossing spans the active-damage phase where the rate is positive). By \ref{L3}, $\mathbb P[\gamma^\dagger]=e^{-\Delta\Sigma(\gamma)}\,\mathbb P[\gamma]$, hence
\[
\mathbb P[\mathcal R]=\sum_{\gamma\in\text{cross}}e^{-\Delta\Sigma(\gamma)}\,\mathbb P[\gamma]\;\le\;\sum_{\text{all }\gamma}e^{-\Delta\Sigma(\gamma)}\,\mathbb P[\gamma]=\E\!\left[e^{-\Delta\Sigma}\right]=1,
\]
the last equality being the integral fluctuation theorem (the detailed relation \ref{L3} summed over all paths) \citep{seifert2012}. The inequality is strict because the non-crossing paths carry positive probability and contribute positively to the full sum but not to $\mathbb P[\mathcal R]$; hence $\mathbb P[\mathcal R]<1$, which is Axiom~\ref{ax:irr}. (The naive bound ``$\Delta\Sigma>0$ with positive probability $\Rightarrow\E[e^{-\Delta\Sigma}]<1$'' is invalid, since reversed paths can carry $\Delta\Sigma<0$; the integral fluctuation theorem is what makes the argument rigorous.) Self-reference (\ref{ax:self}) is what makes $\gamma^\dagger\ne\gamma^{-1}$: because $F_t$ carries the history, the reverse state path is not generated by inverting the forward operator, so the entropy production is genuine rather than a reparametrization artifact.
\end{proof}

\begin{remark}[Nulls as violated regularity]\label{rem:nulls}
Theorem~\ref{thm:reduction} predicts its own failures. A substrate whose recovery balances damage ($\E[\delta_t]\to0$, violating the strict-rise part of \ref{L2}) shows no transition; a symmetric self-reference ($\gamma^\dagger=\gamma^{-1}$, violating \ref{L3}) shows no irreversibility; an externally imposed sequence with no path asymmetry violates \ref{L3} as well. In each case the absence of a critical transition is a confirmed prediction of the reduction, not a counterexample to it.
\end{remark}

% ============================================================
\section{Two-Scale Renormalization}\label{sec:renorm}
% ============================================================

A self-referential system is observed at several scales --- token, turn, sequence, corpus, population --- and the framework requires that the same structural law describe each. We make the inter-scale map precise.

\begin{definition}[Renormalization map]\label{def:M}
For scales $\ell<\ell'$ the renormalization map is the composition
\begin{equation}\label{eq:M}
M_{\ell\to\ell'}=N_{\ell'}\circ R_{\ell\to\ell'}\circ C_{\ell\to\ell'},
\end{equation}
where $C$ \emph{coarse-grains} micro-trajectories into blocks (windows, clusters, representative volumes), $R$ \emph{rescales} the effective parameters along the induced parameter flow $\theta_\ell\mapsto\theta_{\ell'}$, and $N$ \emph{normalizes} dimensionless invariants so that the renormalized description stands structurally on a par with the original.
\end{definition}

\begin{definition}[Normal form and scale-covariance ansatz]\label{prop:forminv}
We take the mean dynamics in the Onsager-type normal form
\begin{equation}\label{eq:normalform}
\dot z = F(z,\Sigma)-\nabla_z\Pi(z,\Sigma)+\Theta\,\nabla_z I(z),
\end{equation}
with drift $F$, dissipation potential $\Pi$, and information coupling $I$: the gradient form of Assumption~\ref{R:flow} in which $-\nabla_z\Pi$ is the dissipation of Axiom~\ref{ax:diss} and $\Theta\,\nabla_z I$ is the self-referential force of Axiom~\ref{ax:self}. We restrict attention to the \emph{relevant sector}: the subspace of operators on which $M_{\ell\to\ell'}$ acts form-covariantly, returning \eqref{eq:normalform} with renormalized $(F',\Pi',\Theta')=R_{\ell\to\ell'}(F,\Pi,\Theta)$.
\end{definition}
\begin{remark}[On the ansatz]
Renormalization generically generates new operators, so exact form-closure is not automatic; the restriction to the relevant sector is the standard content of a renormalization-group analysis, in which irrelevant operators are suppressed under coarse-graining and the surviving form-covariant subspace carries the universal structure \citep{kadanoff1966,wilson1971}. Within this sector the construction is consistent: coarse-graining $C$ replaces $z$ by a block average whose fluctuations contribute, to leading order, a renormalized diffusion absorbed into $\Pi'$ (Assumption~\ref{R:ergodic}), and each gradient term pulls back to a gradient term under the smooth change of variables induced by $R$. Downstream results use only this covariance, not a claim that every microscopic dynamics lies in the relevant sector.
\end{remark}

The five operative scales are: tokens ($\ell_0$), turn-pairs ($\ell_1$), the sequence ($\ell_2$, carrier of the inflection $\hat a$ and the accumulator), the corpus ($\ell_3$, carrier of the $\hat a$-distribution), and the population ($\ell_4$). The scale-covariance of Definition~\ref{prop:forminv} is what licenses transporting the order parameter $\hat a$ across them.

% ============================================================
\section{The Operational Frame}\label{sec:operational}
% ============================================================

The abstract class of Section~\ref{sec:axioms} is instantiated, at the sequence scale $\ell_2$, by a measurement system on embedded trajectories. The instantiation is metric-agnostic in principle; we record the current operational definitions so that the abstract objects $(\Sigma,\hat a)$ have explicit observable counterparts.

\begin{definition}[Deformation observables]\label{def:axes}
Let $(z_t)_{t\le T}$, $z_t=(u_t,a_t)\in\R^{2d}$, be the embedded turn-pair sequence and $v_t:=\norm{z_t-z_{t-1}}$ the semantic velocity. Five coupled deformation axes are defined on windows $\mathcal W(t)$:
\begin{enumerate}[label=\textnormal{(D\arabic*)},leftmargin=3em]
\item \textbf{Stiffness} $E_a:=1/\Var_{\mathcal W}(v)$, inverse velocity variance (strain hardening);
\item \textbf{Asymmetry} $\mathrm{MAS}$, the meta-asymmetry between the user and assistant sub-trajectories in $\R^{d}$ (one-sided load introduction, cantilever);
\item \textbf{Accumulation} $B_w$, the windowed budget accumulator, a monotone functional of the cumulative geometric displacement (creep/fatigue);
\item \textbf{Brittleness} $\mathrm{dir\_consistency}$, the directional coherence of successive increments, whose collapse marks brittle break-out;
\item \textbf{Coupling} $\mathrm{xcorr}(v,w)$, the cross-correlation of semantic velocity $v$ and a structural displacement field $w$ (dynamics--structure covariation).
\end{enumerate}
The order parameter is the \textbf{inflection} $\hat a$: the location of the unique inflection (Theorem~\ref{thm:reduction}(ii)) of the sigmoid fit of cumulative damage $\Sigma$ against cumulative strain.
\end{definition}

\begin{remark}[Anti-sparse coupling]\label{rem:antisparse}
The five axes are coupled, not redundant. In the linear-response (Hooke) regime $\sigma=E\varepsilon$ ties stress, stiffness and strain analytically; the same coupling recurs among the axes, so stiffness (D1) and accumulation (D3) are two projections of one material response. A sparsity-inducing ($\ell_1$) selection would discard such coupled coordinates as collinear and thereby lose information, not noise: the correct paradigm is materials testing, where coupled gauges record one response from distinct axes, not feature selection over independent predictors. Consistently, the loading-time autocorrelation of the cumulative family is signal: measurements along a deformation curve are not independent, and specimen-level inference must use the whole curve as the unit (block resampling), not the individual measurement.
\end{remark}

\begin{definition}[Budget state machine]\label{def:budget}
The accumulator $B_w$ drives a monotone hysteretic state machine with ordered regimes $\textsc{Normal}\preceq\textsc{Watch}\preceq\textsc{Warning}\preceq\textsc{Critical}$, transitions of which are governed by policy thresholds and a recovery rate. By the No-Go theorem (Section~\ref{sec:nogo}) these thresholds are calibration parameters, not universal constants; only the \emph{architecture} (monotone accumulation with hysteretic, asymmetric transitions, Axioms~\ref{ax:acc},\ref{ax:irr}) is part of the theory.
\end{definition}

The embedding instrument is a sentence encoder satisfying Assumption~\ref{R:embed} \citep{reimers2019}; decreasing embedding variance signals stiffening of the semantic space, and irreversibly growing cumulative displacement documents plastic deformation, the observable face of Axioms~\ref{ax:acc} and \ref{ax:irr}.

% ============================================================
\section{T1: Information--Stability Correspondence}\label{sec:t1}
% ============================================================

The first structural theorem links dynamical stability, measured by the maximal Lyapunov exponent, to the trend of mutual information between the coupled coordinates $p:=u$ (prompt/user) and $r:=a$ (response/assistant), $z=(p,r)$.

\begin{lemma}[Covariance growth]\label{lem:covgrowth}
Under Assumptions~\ref{R:flow}--\ref{R:embed}, if $\lambda_{\max}>0$ then the state covariance $\Cov(t)$ satisfies $\norm{\Cov(t)}_{\mathrm{op}}\ge c\,e^{2\lambda_{\max}t}$ for some $c>0$ and all large $t$.
\end{lemma}
\begin{proof}
By Assumption~\ref{R:flow} a perturbation obeys the variational equation $\dot\xi=DF(z(t))\xi$. The covariance solves the Lyapunov equation $\dot\Cov=DF\,\Cov+\Cov\,DF^{\top}+\sigma^2 I$. For the homogeneous part the solution is $\Cov_h(t)=\Phi(t)\Cov(0)\Phi(t)^{\top}$ with fundamental matrix $\Phi$. By Assumption~\ref{R:lyap} and the multiplicative ergodic theorem \citep{oseledets1968}, $\tfrac1t\log\norm{\Phi(t)}_{\mathrm{op}}\to\lambda_{\max}$, so choosing $\Cov(0)$ with a component along the leading Oseledets direction gives $\norm{\Cov_h(t)}_{\mathrm{op}}\ge c\,e^{2\lambda_{\max}t}$. The inhomogeneous noise term contributes $\int_0^t\Phi(t)\Phi(s)^{-1}\sigma^2\Phi(s)^{-\top}\Phi(t)^{\top}ds$, which is positive semidefinite and of order at most $e^{2\lambda_{\max}t}$; it cannot reduce the leading growth. Hence the stated lower bound.
\end{proof}

\begin{lemma}[Two-sided entropy bounds]\label{lem:entropy}
Let $X\in\R^n$ have covariance $\Sigma_X$. Then $h(X)\le\tfrac12\log\!\big((2\pi e)^n\det\Sigma_X\big)$, with equality iff $X$ is Gaussian \citep[Thm.~8.6.5]{coverthomas2006}. If moreover the law of $X$ is log-concave (Assumption~\ref{R:ergodic}), there is a dimensional constant $c_n>0$ with $h(X)\ge\tfrac12\log\det\Sigma_X+c_n$.
\end{lemma}
\begin{proof}
The upper bound is the Gaussian maximum-entropy property. For the lower bound, reduce to the isotropic case by $Y:=\Sigma_X^{-1/2}(X-\E X)$, which is log-concave with identity covariance and satisfies $h(X)=h(Y)+\tfrac12\log\det\Sigma_X$. Isotropic log-concave laws obey a uniform entropy lower bound $h(Y)\ge c_n$ by the variance--entropy comparison from the Brascamp--Lieb inequality and the concentration of log-concave measures \citep{brascamplieb1976,bobkov2003}. Combining gives the claim.
\end{proof}

\begin{theorem}[T1: Information--stability correspondence]\label{thm:t1}
Assume \ref{R:flow}--\ref{R:mi} together with the \emph{coupling non-degeneracy} condition \textnormal{(N)}: in the expanding regime the cross-covariance grows strictly slower than the geometric mean of the marginal blocks, $\norm{\Sigma_{pr}(t)}=o\!\big(\sqrt{\norm{\Sigma_{pp}(t)}\,\norm{\Sigma_{rr}(t)}}\big)$, and in the contracting regime the conditional law of $r$ given $p$ has a bounded noise floor. Then, with $I(t):=I(p_t;r_t)$,
\[
\sgn(\lambda_{\max})=-\sgn\!\Big(\tfrac{d}{dt}I(t)\Big):
\]
instability ($\lambda_{\max}>0$) forces coupling decay, stability ($\lambda_{\max}<0$) forces coupling growth, and criticality ($\lambda_{\max}=0$) is stationary coupling.
\end{theorem}
\begin{proof}
Write $I=h(p)+h(r)-h(p,r)$.

\emph{Case $\lambda_{\max}>0$.} By Lemma~\ref{lem:covgrowth} every covariance block grows like $e^{2\lambda_{\max}t}$, so by Lemma~\ref{lem:entropy} the marginal entropies $h(p),h(r)$ and the joint $h(p,r)$ all grow linearly in $t$ at rate $\propto\lambda_{\max}$. By condition (N) the off-diagonal covariance grows strictly slower than the geometric mean of the diagonal blocks, so the correlation coefficient $\rho_{pr}=\Sigma_{pr}/\sqrt{\Sigma_{pp}\Sigma_{rr}}\to0$. As $\rho_{pr}\to0$ the joint law tends, in the Gaussian/log-concave leading order of Lemma~\ref{lem:entropy}, to the product of the marginals, so $I(p_t;r_t)\to0$; hence $dI/dt<0$.

\emph{Case $\lambda_{\max}<0$.} By Assumption~\ref{R:ergodic} the covariance is bounded (contraction toward the invariant law). Contraction merges nearby trajectories, so the dependence $r=g(p)+\xi$ sharpens toward the bounded noise floor of (N): $h(r\mid p)$ decreases toward the noise-floor entropy while $h(r)$ stays bounded below. Then $I=h(r)-h(r\mid p)$ is nondecreasing, and strictly increasing while the system still approaches stationarity, i.e.\ $dI/dt\ge0$, with $dI/dt>0$ off stationarity.

\emph{Case $\lambda_{\max}=0$.} Neither growth nor contraction dominates; the covariance and the conditional structure are asymptotically constant, so $I$ reaches a stationary value and $dI/dt=0$. This is exactly the boundary at which the sign of $dI/dt$ changes.

Collecting the three cases, $\sgn(\lambda_{\max})=-\sgn(dI/dt)$.
\end{proof}

\begin{remark}[Excluded degeneracy]\label{rem:t1deg}
Condition (N) is necessary, not cosmetic. If a single unstable mode drives both blocks at once (the leading Oseledets direction projecting onto $p$ and $r$ with comparable weight), then $\rho_{pr}\to1$ and instability \emph{increases} the coupling, reversing the correspondence. (N) excludes exactly this synchronization; it holds whenever the two channels are not slaved to one common expanding mode.
\end{remark}

\begin{remark}
T1 is a diagnostic: a declining mutual-information trend signals dynamical instability without computing Lyapunov exponents, and growing or stable coupling signals stability or criticality. It is the abstract content of axis (D5) of Definition~\ref{def:axes}.
\end{remark}

% ============================================================
\section{T2: Autopoietic Stationarity}\label{sec:t2}
% ============================================================

The second theorem gives conditions under which the operator recursion of Definition~\ref{def:srds} has a unique, globally attracting non-equilibrium steady state. It requires one structural hypothesis beyond the axioms.

\begin{definition}[Contraction condition]\label{cond:contraction}
The meta-operator $\Phi$ acts on a complete metric space $(\mathcal F,d_{\mathcal F})$ as a $\kappa$-contraction in its first argument: there is $\kappa\in(0,1)$ with $d_{\mathcal F}(\Phi(F_1,\cdot),\Phi(F_2,\cdot))\le\kappa\,d_{\mathcal F}(F_1,F_2)$ for all $F_1,F_2$. Moreover the stationary state drift is \emph{two-sided dissipative} about its mean $z^*$: there exist $0<c\le c'<\infty$ with $-c'\norm{z-z^*}^2\le\langle z-z^*,F^*(z)\rangle\le-c\,\norm{z-z^*}^2$ on the basin. (This last is a hypothesis on the \emph{state} drift; Axiom~\ref{ax:diss} concerns the operator-level functional and does not by itself bound $F^*$.)
\end{definition}

\begin{theorem}[T2: Autopoietic stationarity]\label{thm:t2}
Under Axioms~\ref{ax:self}--\ref{ax:diss} and Definition~\ref{cond:contraction}:
\begin{enumerate}[label=\textnormal{(\roman*)},leftmargin=2.2em]
\item There is a unique autopoietic fixed point $F^*$ with $\Phi(F^*)=F^*$, and the iterates $F_{n+1}=\Phi(F_n)$ converge geometrically, $d_{\mathcal F}(F_n,F^*)\le\frac{\kappa^{n}}{1-\kappa}\,d_{\mathcal F}(F_1,F_0)$.
\item The state dynamics under $F^*$ possess a globally attracting non-equilibrium regime whose stationary fluctuation level is bounded away from zero, in the asymptotically absorbing band $[\,\sigma^2 m/(4c'),\,\sigma^2 m/(4c)\,]$.
\item At $F^*$ the dissipation functional is stationary, $\dot D=0$.
\end{enumerate}
\end{theorem}
\begin{proof}
\textbf{(i)} $(\mathcal F,d_{\mathcal F})$ is complete and $\Phi$ is a $\kappa$-contraction (Definition~\ref{cond:contraction}); the Banach fixed-point theorem \citep{banach1922} yields a unique $F^*$ with $\Phi(F^*)=F^*$ and, for any $F_0$, the a-priori estimate $d_{\mathcal F}(F_n,F^*)\le\frac{\kappa^n}{1-\kappa}d_{\mathcal F}(F_1,F_0)$.

\textbf{(ii)} Fix the operator at $F^*$; the state then solves the autonomous diffusion $dz=F^*(z)\,dt+\sigma\,dW_t$ (Assumption~\ref{R:flow}) with stationary mean $z^*=\E_\mu[z]$. Define the quadratic Lyapunov functional $\mathcal V(z):=\tfrac12\norm{z-z^*}^2$. By It\^o's formula,
\[
d\mathcal V=\big\langle z-z^*,\,F^*(z)\big\rangle\,dt+\tfrac12\sigma^2 m\,dt+\sigma\langle z-z^*,dW_t\rangle .
\]
The two-sided dissipativity of Definition~\ref{cond:contraction} gives $-2c'\,\mathcal V(z)\le\langle z-z^*,F^*(z)\rangle\le-2c\,\mathcal V(z)$. Write $\bar V(t):=\E[\mathcal V(z(t))]$; taking expectations and using the It\^o noise injection rate $\tfrac12\sigma^2 m$, where $m$ is the state dimension of Assumption~\ref{R:flow}, the two bounds sandwich the mean Lyapunov rate:
\[
-2c'\,\bar V+\tfrac12\sigma^2 m\;\le\;\dot{\bar V}\;\le\;-2c\,\bar V+\tfrac12\sigma^2 m .
\]
Both scalar comparison ODEs are globally attracting (to $\sigma^2 m/(4c')$ and $\sigma^2 m/(4c)$ respectively), so by comparison
\[
\liminf_{t\to\infty}\bar V(t)\ge\tfrac{\sigma^2 m}{4c'}>0,\qquad \limsup_{t\to\infty}\bar V(t)\le\tfrac{\sigma^2 m}{4c}.
\]
The band $[\sigma^2 m/(4c'),\sigma^2 m/(4c)]$ is asymptotically absorbing (the limit exists when $c=c'$). The $\liminf$ bound keeps the stationary fluctuation level bounded away from $0$ (the system maintains finite fluctuations rather than collapsing to a deterministic point, non-equilibrium); absorption into the band is global attraction, stability.

\textbf{(iii)} Let $D:=\E[C(z,u)]-C_{\mathrm{eq}}$ be the dissipation functional, with $C$ the cost whose minimization is the gradient structure of \eqref{eq:normalform}. Then $\dot D=\partial_t\E[C]$. In the stationary regime of (ii) the law of $z$ is time-invariant, so $\partial_t\E[C]=0$; equivalently $z^*$ is a critical point of $C$ along $F^*$, $\langle\nabla_z C(z^*),F^*(z^*)\rangle=0$. Hence $\dot D=0$: entropy production is exactly balanced by dissipative export, the defining balance of a non-equilibrium steady state \citep{prigogine1977}.
\end{proof}

\begin{remark}
T2 is the existence statement behind the autopoietic configuration of Definition~\ref{def:srds}: self-reference plus dissipation plus contraction produces a stable far-from-equilibrium attractor, not a dead equilibrium. The geometric rate in (i) is conditional on a \emph{uniform} modulus $\kappa<1$; if the per-step contraction degrades ($\kappa_n\to1$), convergence is governed by the cumulative budget rather than a fixed rate (Section~\ref{sec:spectral}).
\end{remark}

% ============================================================
\section{T3: The Goldilocks Principle}\label{sec:t3}
% ============================================================

The third result concerns where, relative to criticality, a self-referential system functions best. It is stated as a hypothesis, with a proved partial result; the hypothesis status is intrinsic, because the performance functional is not fixed by the axioms.

\begin{hypothesis}[Goldilocks principle]\label{hyp:goldi}
Let $L(\mathrm{I},\mathrm{C})$ be a performance functional in innovation $\mathrm{I}$ (exploration/variety) and coherence $\mathrm{C}$ (stability/recurrence), jointly concave with positive cross-derivative $\partial^2 L/\partial\mathrm{I}\,\partial\mathrm{C}>0$. Then optimal performance requires interior balance, neither $\mathrm{I}=0$ nor $\mathrm{C}=0$, and the optimum lies in a neighbourhood of criticality $\abs{\lambda_{\max}}\approx0$ rather than at the measure-zero point of exact criticality.
\end{hypothesis}

\begin{proposition}[Necessity of interior balance]\label{prop:goldi}
Let $L:[0,\infty)^2\to\R$ be $C^2$ and strictly concave, with positive cross-partial $\partial^2 L/\partial\mathrm I\,\partial\mathrm C>0$ and strictly positive boundary gradients $\partial_{\mathrm I}L(0,0)>0$, $\partial_{\mathrm C}L(0,0)>0$. If $L$ attains a maximum on $[0,\infty)^2$, it is at a unique interior point $(\mathrm I^*,\mathrm C^*)$ with $\mathrm I^*>0$, $\mathrm C^*>0$, and every boundary point is strictly suboptimal.
\end{proposition}
\begin{proof}
A strictly concave function on a convex domain has at most one local maximum, which is then global. At an interior maximum the first-order conditions $\partial_{\mathrm I}L=\partial_{\mathrm C}L=0$ hold. Suppose the maximum were on the boundary, say at $\mathrm I=0$; optimality on that face requires $\partial_{\mathrm I}L\le0$ along $\mathrm I=0$. But by the positive cross-partial and the strict boundary-gradient hypotheses, for every $\mathrm C\ge0$,
\[
\partial_{\mathrm I}L(0,\mathrm C)=\partial_{\mathrm I}L(0,0)+\int_0^{\mathrm C}\partial^2_{\mathrm I\mathrm C}L\,d\mathrm C'\;\ge\;\partial_{\mathrm I}L(0,0)\;>\;0,
\]
since the integrand is positive (the integral is $\ge0$, vanishing only at the corner $\mathrm C=0$). This contradicts boundary optimality on the entire face $\mathrm I=0$, the corner $(0,0)$ included; the maximum is therefore interior with $\mathrm I^*>0$, and symmetrically $\mathrm C^*>0$. Strict concavity makes the interior critical point the unique global maximum, with every boundary point strictly below it.
\end{proof}

\begin{remark}[Why a hypothesis]
Proposition~\ref{prop:goldi} proves the interior-balance \emph{structure}; the full Hypothesis~\ref{hyp:goldi} remains open because (a) the functional $L$ is application-dependent and not derivable from the axioms, and (b) the identification of the interior optimum with the neighbourhood of $\lambda_{\max}=0$ is an empirical prediction: systems operating near, but not at, criticality outperform both rigid ($\mathrm C$ saturated, $\lambda_{\max}\ll0$) and chaotic ($\lambda_{\max}\gg0$) regimes. The principle is the optimization counterpart of T1: criticality is the boundary where the information--stability sign flips, and function concentrates in its neighbourhood.
\end{remark}

% ============================================================
\section{The Renormalization Operator and the Inflection Fixed Point}\label{sec:operator}
% ============================================================

We now realize the inflection $\hat a$ of Axiom~\ref{ax:trans} as the order parameter of a concrete operator fixed point, and quantify its convergence.

\begin{definition}[Damped renormalization operator]\label{def:Reta}
On the cone $\mathcal K\subset C([0,1])$ of strictly positive profiles, with anchored monotone profiles $\mathcal M$ as in Section~\ref{sec:prelim}, let $T_\varepsilon$ be smoothing by a strictly positive, totally positive of order two (TP2) kernel at scale $\varepsilon$, and let
\[
\Rc_{\eta,\varepsilon}:=\Ac\circ\big((1-\eta)\,T_\varepsilon+\eta\,\Pi_{\mathrm{ref}}\big),
\]
where $\Pi_{\mathrm{ref}}$ is damping toward a reference profile with weight $\eta\in[0,1)$ and $\Ac$ is an anchoring gauge that fixes the boundary normalization. Equip $\mathcal K$ with the Hilbert projective metric $d_H$.
\end{definition}

\begin{theorem}[Existence and uniqueness of the monotone fixed point]\label{thm:fixedpoint}
For each fixed $\varepsilon>0$, $\Rc_{\eta,\varepsilon}$ has a unique fixed point $Q^*\in\mathcal M$, and $d_H(\Rc^n f,Q^*)\le\kappa(\varepsilon)^n\,d_H(f,Q^*)$ for all $f\in\mathcal K$, with $\kappa(\varepsilon)=\tanh(\Delta(T_\varepsilon)/4)<1$.
\end{theorem}
\begin{proof}
A strictly positive kernel has finite projective diameter $\Delta(T_\varepsilon)<\infty$, so by Birkhoff's contraction theorem \citep{birkhoff1957,bushell1973,evesonnussbaum1995} $T_\varepsilon$ contracts $d_H$ with coefficient $\tanh(\Delta/4)<1$; the convex damping by $\Pi_{\mathrm{ref}}$ and the gauge $\Ac$ do not increase $d_H$ (the gauge selects a representative within each projective class and $d_H$ is scale-invariant). Hence $\Rc_{\eta,\varepsilon}$ is a $d_H$-contraction on the complete cone, and Banach's theorem \citep{banach1922} gives the unique fixed point $Q^*$ and the geometric estimate. TP2 preservation keeps the iterates monotone, so $Q^*\in\mathcal M$ \citep{karlin1968}.
\end{proof}

\begin{theorem}[Sharp exponential decay of the global contraction]\label{thm:sharpdecay}
For the Gaussian kernel, $\Delta(T_\varepsilon)=2/\varepsilon$, so
\[
\kappa(\varepsilon)=\tanh\!\Big(\tfrac{1}{2\varepsilon}\Big)=1-2e^{-1/\varepsilon}+O(e^{-2/\varepsilon}),
\qquad 1-\kappa(\varepsilon)\sim 2e^{-1/\varepsilon},
\]
whereas the local diffusive spectral gap closes only polynomially, $1-\lambda_1(\varepsilon)\sim C\varepsilon$.
\end{theorem}
\begin{proof}
The projective diameter of the Gaussian smoothing on monotone profiles is $\Delta(T_\varepsilon)=2/\varepsilon$ (the log-ratio spread of the kernel over the unit interval), so by Theorem~\ref{thm:fixedpoint} $\kappa(\varepsilon)=\tanh(1/(2\varepsilon))$. Expanding $\tanh(z)=\frac{1-e^{-2z}}{1+e^{-2z}}=1-2e^{-2z}+O(e^{-4z})$ at $z=1/(2\varepsilon)$ gives $\kappa(\varepsilon)=1-2e^{-1/\varepsilon}+O(e^{-2/\varepsilon})$. The periodic heat semigroup (standard normalization) has eigenvalues $e^{-4\pi^2 n^2\varepsilon}$, so the local gap is $1-e^{-4\pi^2\varepsilon}\sim 4\pi^2\varepsilon$, polynomial in $\varepsilon$. A change of scale convention rescales $\varepsilon$ but leaves the qualitative contrast intact: global projective synchronization is exponentially harder than local smoothing.
\end{proof}

\begin{theorem}[Critical cooling law]\label{thm:cooling}
For the annealed iteration $Q_{n+1}=\Rc_{\varepsilon_n}(Q_n)$ with $\varepsilon_n=c/\ln n$, the cumulative contraction budget diverges iff $c\ge1$:
\[
\sum_{n\ge2}\big(1-\kappa(\varepsilon_n)\big)=\infty \iff c\ge1 .
\]
\end{theorem}
\begin{proof}
By Theorem~\ref{thm:sharpdecay}, $1-\kappa(\varepsilon_n)\sim 2e^{-1/\varepsilon_n}=2e^{-\ln n/c}=2\,n^{-1/c}$. The series $\sum_n n^{-1/c}$ diverges iff the exponent $1/c\le1$, i.e.\ $c\ge1$ ($p$-series criterion).
\end{proof}

\begin{definition}[Inflection order parameter]\label{def:ahat}
The inflection of the fixed point is $\hat a:=\argmax_{x\in(0,1)}(Q^*)'(x)$, the location of maximal slope of $Q^*$: the operator realization of the transition order parameter of Axiom~\ref{ax:trans}.
\end{definition}

\paragraph{Aggregation modes.} Under loading, the micro-imperfection cascade activates in one of four modes, which fix the form class of $Q^*$: (i) parallel-instantaneous $\to$ logistic (central-limit aggregation); (ii) serial-cascaded $\to$ Weibull (weakest-link order statistics); (iii) multiaxial $\to$ hybrid (Hill/Gompertz); (iv) continuous-accumulating $\to$ Gaussian/logistic (diffusive averaging). A substrate is in general a convex combination on the simplex $\boldsymbol\pi=(\pi_1,\pi_2,\pi_3,\pi_4)\in\Delta_4$; the mode identities are the basis families of the form-class vector space, and neither Theorem~\ref{thm:fixedpoint} nor the No-Go theorem below is affected by the mixture: only the inflection location and the mixture vector carry the substrate-specific asymmetry. The classical ductile/brittle dichotomy is the binary projection of this finer partition.

% ============================================================
\section{The No-Go Theorem}\label{sec:nogo}
% ============================================================

The central negative result: the inflection location, although well defined for each fixed point, is \emph{not} a universal constant across the operator class. This is what makes substrate-specific asymmetry a measurement task rather than a defect.

\begin{definition}[Reparametrization action and equivariant transition functional]
The increasing homeomorphisms $h:[0,1]\to[0,1]$ act on profiles by $(h\cdot Q)(x):=Q(h(x))$, with $h^{-1}\cdot(h\cdot Q)=Q$. A class $\mathcal G$ of renormalization operators (TP2, strictly positive, anchored, damped) is \emph{reparametrization-closed} if $R\in\mathcal G$ implies the conjugate $R^{(h)}:=h^{-1}\circ R\circ h\in\mathcal G$ for every such $h$. A transition-location functional $\tau:\mathcal M\to(0,1)$ is \emph{equivariant} if $\tau(h\cdot Q)=h^{-1}(\tau(Q))$.
\end{definition}

\begin{theorem}[No-Go: no universal inflection constant]\label{thm:nogo}
Let $\mathcal G$ be reparametrization-closed and $\tau$ equivariant. Then there is no constant $c\in(0,1)$ with $\tau(Q^*)=c$ for the fixed points $Q^*$ of all $R\in\mathcal G$.
\end{theorem}
\begin{proof}
Suppose such a $c$ exists. Fix $R\in\mathcal G$ with fixed point $Q^*$, so $\tau(Q^*)=c$. Let $\tilde c\in(0,1)$ be arbitrary and choose an increasing homeomorphism $h$ with $h(c)=\tilde c$ (any smooth monotone $h$ with $h(0)=0,h(1)=1,h(c)=\tilde c$ works). Form the conjugate $R^{(h)}=h^{-1}\circ R\circ h\in\mathcal G$ (reparametrization closure) and the profile $\tilde Q^*:=h^{-1}\cdot Q^*$. Then
\[
R^{(h)}(\tilde Q^*)=(h^{-1}\circ R\circ h)(h^{-1}\cdot Q^*)=h^{-1}\cdot R\big(h\cdot(h^{-1}\cdot Q^*)\big)=h^{-1}\cdot R(Q^*)=h^{-1}\cdot Q^*=\tilde Q^*,
\]
so $\tilde Q^*$ is the fixed point of $R^{(h)}$. By equivariance, $\tau(\tilde Q^*)=\tau(h^{-1}\cdot Q^*)=h(\tau(Q^*))=h(c)=\tilde c$. Since $\tilde c$ was arbitrary, the transition location ranges over all of $(0,1)$ across $\mathcal G$, contradicting $\tau\equiv c$.
\end{proof}

\begin{corollary}[Symmetry recovers the constant]\label{cor:symmetry}
A universal inflection location can hold only under a structural restriction that breaks reparametrization freedom. In particular, if the kernel of $R$ is mirror-symmetric, the fixed-point profile satisfies the central symmetry $Q^*(x)=1-Q^*(1-x)$, and its inflection (Definition~\ref{def:ahat}) is at $1/2$.
\end{corollary}
\begin{proof}
Let $S$ be the central involution $(SQ)(x):=1-Q(1-x)$. It maps $\mathcal M$ to itself: $SQ$ is nondecreasing with $(SQ)(0)=1-Q(1)=0$ and $(SQ)(1)=1-Q(0)=1$ (note that the bare domain reflection $Q\mapsto Q(1-\cdot)$ does \emph{not} preserve $\mathcal M$, since it reverses monotonicity; the range flip restores it). A mirror-symmetric kernel makes $\Rc$ commute with $S$, $\Rc\circ S=S\circ\Rc$; by uniqueness of the fixed point (Theorem~\ref{thm:fixedpoint}), $Q^*=SQ^*$, i.e.\ $Q^*(x)=1-Q^*(1-x)$. Differentiating gives $(Q^*)'(x)=(Q^*)'(1-x)$, so the density is symmetric about $1/2$ and its unique maximum, the inflection of $Q^*$, lies at $x=1/2$. The restriction to the centrally symmetric subclass is exactly the loss of reparametrization freedom required by Theorem~\ref{thm:nogo}.
\end{proof}

\begin{theorem}[Certified non-degeneracy under asymmetry]\label{thm:certified}
For the damped Gaussian family parametrized by kernel asymmetry $\alpha$, the inflection map $\alpha\mapsto\tau(Q^*_\alpha)$ is $C^1$, and its first-order response at the symmetric point does not vanish. Writing
\[
\dot Q(1/2):=\partial_\alpha\,\Xi(\alpha)\big|_{\alpha=0}\in[1.8412,\,1.8963]\neq0
\]
for the \emph{perturbation derivative} of the symmetry-breaking functional $\Xi$ --- the first variation, at the midpoint, of the fixed-point condition under kernel asymmetry --- this is a rigorous enclosure obtained by verified interval arithmetic at $256$-bit precision (Krawczyk operator with Anselone bound and box enclosure) \citep{johansson2017,rump1999}. Consequently $d\tau/d\alpha|_{0}\neq0$: the inflection shifts at first order off the symmetric point. The certified spectral bound $\rho(D\Rc)<0.45$ guarantees geometric convergence at the certified parameters.
\end{theorem}
\begin{proof}
By the uniform contraction of Theorem~\ref{thm:fixedpoint} and the certified bound $\rho(D\Rc)<0.45<1$, the linearization $I-D\Rc$ is invertible, so the implicit function theorem applied to $\Rc_{\alpha}(Q^*_\alpha)=Q^*_\alpha$ yields a $C^1$ branch $\alpha\mapsto Q^*_\alpha$, hence a $C^1$ map $\alpha\mapsto\tau(Q^*_\alpha)$. The first-order shift of the inflection is governed by the linearized inflection condition, whose coefficient is the perturbation derivative $\dot Q(1/2)=\partial_\alpha\Xi(\alpha)|_0$ (not the spatial slope of the profile); its enclosure $[1.8412,1.8963]$ is produced by a reproducible ball-arithmetic certificate with auditable load-bearing steps. Since the enclosure excludes $0$, $d\tau/d\alpha|_0\neq0$, so no neighbourhood of the symmetric kernel shares the inflection value $1/2$. (The inflection is non-degenerate: the symmetric fixed point has a strict density maximum, so the linearized inflection condition has nonzero leading coefficient $(Q^*)'''(1/2)\neq0$ and $d\tau/d\alpha|_0=-\dot Q(1/2)/(Q^*)'''(1/2)$ is well defined.)
\end{proof}

\begin{remark}[Scope]
Theorem~\ref{thm:nogo} is a non-existence statement and is complete as proved. Theorem~\ref{thm:certified} certifies one explicit asymmetric direction; the general statement that $\tau$ shifts for \emph{every} admissible asymmetry, at all $(\varepsilon,\eta)$, remains an open conjecture, but the No-Go theorem does not require it: a single asymmetric kernel at which the inflection moves already refutes universality. The operative order parameter $\hat a=\argmax_x(Q^*)'(x)$ (Definition~\ref{def:ahat}) is itself a location functional: the inflection of a conjugate fixed point $Q^*\circ h$ is $\argmax_x\big((Q^*)'(h(x))\,h'(x)\big)$, which can be placed anywhere in $(0,1)$ by choice of the increasing homeomorphism $h$, so no universal $\hat a$-constant exists either; the abstract No-Go transfers to the concrete order parameter.
\end{remark}

\begin{conjecture}[The self-similar candidate $\Sigma_c$]\label{conj:sigmac}
A minimal self-similarity condition, whole-to-core as core-to-remainder, selects a candidate inflection: $\frac{1}{\Sigma}=\frac{\Sigma}{1-\Sigma}\iff\Sigma^2+\Sigma-1=0$, with positive root $\Sigma_c=\varphi-1=1/\varphi\approx0.618$. Three independent routes (the self-similarity condition, the damage fixed-point iteration $\Sigma_{t+1}=1-\Sigma_t^2$, and a dimensional optimization) converge on this root. This is a conjecture, not a theorem: by Theorem~\ref{thm:nogo} no such value can be universal across $\mathcal G$, and the empirically observed inflections scatter in a structured, substrate-specific way around it. Proximity of an individual measured inflection to $1/\varphi$ is a datum, not support for universality.
\end{conjecture}

% ============================================================
\section[The Second No-Go: No Static Classifier]{The Second No-Go: No Universal Static Classifier under Interventional Monitoring}\label{sec:nogo2}
% ============================================================

Theorem~\ref{thm:nogo} eliminates universal \emph{numerical} structure: no inflection constant survives reparametrization freedom. This section eliminates, at the deployment level of Section~\ref{sec:operational}, universal \emph{verdict} structure: no static binary classifier survives the closure of the measurement--action loop, once the published score enters the subsequent dynamics and the reception is adversive. The two results are projections of one limit picture: what is universal is the structural grammar and the monitoring architecture, never a constant and never a static verdict.

\begin{definition}[Monitoring modes]\label{def:modes}
Let $(x_t)$ be the monitored trajectory and $\hat a_t$ the measured order parameter (Definition~\ref{def:axes}). Monitoring is \emph{observational} if the dynamics do not depend on the measurement,
\[
x_{t+1}=H_0(x_t),
\]
and \emph{interventional} if the published verdict enters the update,
\[
x_{t+1}=H_I\big(x_t,\Phi(\hat a_t)\big),
\]
for a classifier $\Phi$ with values in $\{0,1\}$ ($1=$ \textsc{unstable}). In materials-testing terms the observational mode is non-destructive testing, the specimen is invariant under the measurement, while the interventional mode changes the boundary conditions of the loaded system through the act of reporting.
\end{definition}

\begin{definition}[Self-defeating reception]\label{def:sdelta}
A system $a_0$ in a class $\mathcal A$ has \emph{self-defeating reception} if its realized stability outcome negates every published verdict: writing $C(a_0)\in\{0,1\}$ for the outcome under interventional monitoring,
\[
C(a_0)=1-\Phi(a_0)\qquad\text{for every static }\Phi.
\]
The verdict \textsc{unstable} induces correction (the system stabilizes); the verdict \textsc{stable} induces complacency (the system destabilizes). This is the self-defeating prophecy \citep{merton1948}, recorded here as a class property, not constructed from the axioms. We write
\[
S_\Delta:=\{a\in\mathcal A:\ a\ \text{has self-defeating reception}\}\subseteq\mathcal A
\]
for the subset of such systems.
\end{definition}

\begin{theorem}[Diagonal feedback barrier]\label{thm:nogo2}
Let $\mathcal A$ be a system class with $S_\Delta\ne\emptyset$. Then no static classifier $\Phi:\mathcal A\to\{0,1\}$ is correct on $\mathcal A$ under interventional monitoring.
\end{theorem}
\begin{proof}
Correctness requires $\Phi(a)=C(a)$ for all $a\in\mathcal A$. For $a_0\in S_\Delta$ this reads $\Phi(a_0)=1-\Phi(a_0)$, i.e.\ a fixed point of the negation on $\{0,1\}$; the negation has none.
\end{proof}

\begin{remark}[Status of the theorem]\label{rem:status2}
Theorem~\ref{thm:nogo2} is the terminal step of the classical diagonal argument \citep{lawvere1969,yanofsky2003}: fixed-point freeness of negation. It is not new mathematics, and the negation is \emph{defined into} the class by Definition~\ref{def:sdelta}, not constructed. Its content is locative, parallel to Theorem~\ref{thm:nogo}: it fixes \emph{where} the impossibility lives --- in the closed, adversive measurement--action loop --- and thereby what monitoring may still claim there. Whether a given real class (conversational trajectories under visible scoring among them) actually has $S_\Delta\ne\emptyset$ is an empirical premise: it can be supported by a controlled mode switch $H_0\leftrightarrow H_I$, never settled formally.
\end{remark}

\begin{proposition}[A closed loop is necessary, not sufficient]\label{prop:loop}
Let the reception of the published verdict be $\psi:\{0,1\}\to\{0,1\}$, so that the realized outcome at the loop is $C=\psi\circ\Phi$. Among the four self-maps of $\{0,1\}$, the negation is the only fixed-point-free one. Consequently: if $\psi\ne\lnot$ (in particular $\psi=\mathrm{id}$, the self-fulfilling reception, and the constant receptions), then $\psi$ has a fixed point $b^*$ and the constant classifier $\Phi\equiv b^*$ is correct at the loop: a regime-bound classifier exists. Only the adversive loop $\psi=\lnot$ raises the barrier.
\end{proposition}
\begin{proof}
$\mathrm{id}$ fixes both points, the constants fix their value, and $\lnot$ exchanges the two points. If $\psi(b^*)=b^*$, publishing $b^*$ yields outcome $\psi(b^*)=b^*$, which is the published verdict.
\end{proof}

\begin{remark}[The $\Phi_{\mathcal R}$ diagnosis: two edges of one fixed-point picture]\label{rem:phir}
The existence question for a correct classifier is in both classical traditions a fixed-point question for the reaction map, answered at two different edges. On the continuous edge, a public prediction with \emph{continuous} reaction admits a correct value by the intermediate-value/Brouwer argument, the classical possibility theorem of public prediction \citep{grunbergmodigliani1954,simon1954}. On the binary edge, Proposition~\ref{prop:loop} locates the unique exception: only the adversive reception is fixed-point free. Continuity buys the fixed point; discreteness localizes its only failure.

The tempting strengthening, ``a regime-bound classifier exists \emph{iff} no reception coordinate is the negation'', holds only under \emph{separability} (memoryless receptions acting coordinatewise, $\psi=(r_a)_{a}$). It fails in the coupled case: on $\{0,1\}^2$ the rotation $\psi(x,y)=(y,\,1{-}x)$ is fixed-point free although neither coordinate negates itself: of the nine fixed-point-free permutations of the four states, only three are separable. Adversiveness can be \emph{distributed} across coupled axes ($x$ follows $y$, $y$ opposes $x$) without being attributable to any single axis. Since specimens carry loading-path memory --- $\hat a_t$ is a functional of the whole history $x_{0:t}$ (Axioms~\ref{ax:acc}, \ref{ax:irr}), and the deformation axes are coupled (Remark~\ref{rem:antisparse}) --- the non-separable case is the generic one for this framework. The characterization therefore remains a \emph{diagnosis}, not a theorem: it unifies the two edges and names the exact residue (the sign of the coupled loop), which is an empirical question, not a structural one.

The design consequence is the operational counterpart of Theorem~\ref{thm:nogo}: monitoring must be specified as $\Phi_{\mathcal R}$ --- bound to an explicit observation, feedback and intervention regime $\mathcal R$ --- and never as a universal $\Phi_{\mathrm{univ}}$. This is the architecture already adopted in Definition~\ref{def:budget}: thresholds are calibration parameters of a regime, the state machine is the invariant. The continuous-optimization literature on prediction that shifts its own distribution \citep{perdomo2020} studies retraining equilibria for risk minimizers; the binary verdict cut and its barrier are a separate, elementary stratum beneath it.
\end{remark}

\begin{remark}[Two No-Gos, one picture]\label{rem:twonogos}
Theorem~\ref{thm:nogo} forbids a universal \emph{value} by reparametrization freedom of the substrate; Theorem~\ref{thm:nogo2} forbids a universal \emph{static verdict} by fixed-point freeness of the adversive loop. Both eliminate universals in favour of structure: the functional form, the accumulation architecture, and the regime-bound measurement survive; constants and context-free classifications do not. The first No-Go makes the inflection a measured quantity; the second makes the verdict a regime property. Together they state the honest scope of the framework's operational claim: structural grammar universal, numbers and verdicts local.
\end{remark}

% ============================================================
\section{The Elimination Ledger}\label{sec:ledger}
% ============================================================

We collect, in reader-checkable form, what \emph{cannot} be concluded at the present level of generality, and the minimal structure that would change each conclusion. These are not statements of empirical uncertainty: each is a consequence of the preceding theorems.

\begin{proposition}[E1: no universal rate from local contractivity]\label{E1}
Per-step contractivity (existence of $\kappa_n<1$) guarantees convergence once the budget diverges (Theorem~\ref{thm:cooling}), but does not determine a convergence-rate class (exponential / polynomial / \dots) across systems: the rate is set by the cumulative budget $A_n:=\sum_{k\le n}(1-\kappa_k)$, not by any single $\kappa_n$.
\end{proposition}
\begin{proof}
By Theorem~\ref{thm:fixedpoint}, $d_H(Q_n,Q^*)\le\big(\prod_{k\le n}\kappa_k\big)d_H(Q_0,Q^*)\le e^{-A_n}d_H(Q_0,Q^*)$, using $\kappa_k\le e^{-(1-\kappa_k)}$. The bound depends on $A_n$ only; the example $\kappa_n=\tanh(1/(2\varepsilon_n))$ with $\varepsilon_n=c/\ln n$ (Theorem~\ref{thm:cooling}) gives $A_n\sim\sum n^{-1/c}$, whose growth class (and hence the rate) is tuned continuously by $c$ while every local $\kappa_n<1$.
\end{proof}

\begin{proposition}[E2: sharpness of budget bounds]\label{E2}
Even with an a-priori bound $d_n\le e^{-A_n}d_0$, no universal improvement holds within the class: for any prescribed decay profile there is an admissible schedule $(\varepsilon_n)$ realizing it.
\end{proposition}
\begin{proof}
Given a target nonincreasing profile $\beta_n\downarrow0$, choose $\kappa_n$ with $\prod_{k\le n}\kappa_k=\beta_n$ (i.e.\ $1-\kappa_n=1-\beta_n/\beta_{n-1}$), realizable by $\varepsilon_n=1/\big(2\,\mathrm{artanh}\,\kappa_n\big)$ since $\kappa\mapsto\tanh(1/(2\varepsilon))$ is a bijection onto $(0,1)$. Each $\kappa_n<1$ is admissible, so the decay can be made arbitrarily slow.
\end{proof}

\begin{proposition}[E3: schedule dependence]\label{E3}
When the per-step contraction degrades ($\kappa_n\to1$, e.g.\ $\varepsilon_n\to0$), convergence hinges on the explicit divergence $\sum_n(1-\kappa_n)=\infty$ (Theorem~\ref{thm:cooling}); absent this schedule condition, persistent bad windows or non-convergence cannot be excluded.
\end{proposition}

\begin{proposition}[E4: no universal numerical threshold]\label{E4}
No concrete inflection/transition value (e.g.\ $1/2$, $1/\varphi$) is derivable at this level of generality; by Theorem~\ref{thm:nogo} any such value requires structure that breaks reparametrization freedom.
\end{proposition}

\begin{proposition}[E5: no universal static classifier]\label{E5}
Correctness of a static classifier $\Phi:\mathcal A\to\{0,1\}$ on a feedback-responsive class cannot be concluded without an empirically established reception sign for the deployed regime: if $S_\Delta\ne\emptyset$, no such $\Phi$ exists (Theorem~\ref{thm:nogo2}); only $\Phi_{\mathcal R}$-statements relative to a specified observation--feedback--intervention regime are admissible.
\end{proposition}

\paragraph{Minimal structure that changes the conclusions.}
\begin{description}[leftmargin=2.4em,style=nextline]
\item[\textnormal{(S1) Uniform minorization (Doeblin).}] A genuine lower bound $Tf\ge c\,\ip{f}{\nu}$ yields uniform mixing and a uniform rate (Section~\ref{sec:spectral}); converts E1/E3 into an explicit exponential rate.
\item[\textnormal{(S2) Time-structure (good-step density / bounded gaps).}] A lower bound on the frequency of contractive steps converts local contractivity into an explicit rate in terms of that density; addresses E1/E2.
\item[\textnormal{(S3) Symmetry / identifiability.}] Restriction to a mirror-symmetric or identifiable normal-form subclass pins a numerical inflection location (Corollary~\ref{cor:symmetry}); addresses E4.
\end{description}

% ============================================================
\section{Spectral and Contraction Substructure}\label{sec:spectral}
% ============================================================

The rate of approach to the autopoietic fixed point is governed by the contraction budget, not by any local coefficient. We make the governing structure precise and exhibit the uniform-rate restoration (S1).

\begin{definition}[Contraction budget]
For a (possibly non-autonomous) iteration with per-step coefficients $\kappa_n\in(0,1)$, the cumulative contraction budget is $A_n:=\sum_{k\le n}(1-\kappa_k)$.
\end{definition}

\begin{theorem}[Budget governs convergence]\label{thm:budget}
$d_H(Q_n,Q^*)\le e^{-A_n}d_H(Q_0,Q^*)$; in particular $A_n\to\infty$ implies $Q_n\to Q^*$, while a bounded budget leaves a non-vanishing contraction factor and so cannot force convergence. The asymptotic rate class is determined by the growth class of $A_n$ and is invariant under any reparametrization that preserves $A_n$.
\end{theorem}
\begin{proof}
The product bound $\prod_{k\le n}\kappa_k\le e^{-\sum_{k\le n}(1-\kappa_k)}=e^{-A_n}$ gives the estimate; if $A_n\to\infty$ then $d_H(Q_n,Q^*)\to0$. Conversely, if $A_n$ is bounded then $\prod\kappa_k\ge\prod e^{-2(1-\kappa_k)}=e^{-2A_\infty}>0$ (using $\kappa\ge e^{-2(1-\kappa)}$ for $\kappa$ near $1$), so the contraction factor does not vanish and convergence to a unique limit cannot be forced. Reparametrization-invariance of the rate follows because $A_n$ is built from the projective coefficients $\kappa_k$, which are gauge-invariant (Theorem~\ref{thm:fixedpoint}).
\end{proof}

\begin{theorem}[Doeblin restart: uniform rate (S1)]\label{thm:doeblin}
Replace $T_\varepsilon$ by the regularized kernel $k_{\eta,\varepsilon}=(1-\eta)k_\varepsilon+\eta$ (a uniform minorization of weight $\eta$). Then the projective contraction satisfies $\kappa_\eta\le\frac{1-\eta}{1+\eta}<1$ \emph{uniformly in} $\varepsilon$, and the annealed iteration converges provided $\sum_n\eta_n=\infty$, independently of the cooling schedule $\varepsilon_n$.
\end{theorem}
\begin{proof}
A uniform additive component $\eta$ is a Doeblin minorization $k_{\eta,\varepsilon}\ge\eta$; by the Doeblin coupling bound \citep{doeblin1937,seneta2006} the Hilbert/total-variation contraction is at most $\frac{1-\eta}{1+\eta}$, with no dependence on $\varepsilon$. Hence $1-\kappa_{\eta_n}\ge\frac{2\eta_n}{1+\eta_n}\ge\eta_n$, so the budget obeys $A_n\ge\sum_{k\le n}\eta_k$; by Theorem~\ref{thm:budget}, $\sum_n\eta_n=\infty$ forces convergence regardless of $\varepsilon_n$.
\end{proof}

\begin{remark}[Local versus global gap]
Theorem~\ref{thm:sharpdecay} exhibits the central discrepancy: the \emph{local} diffusive relaxation gap closes polynomially ($1-\lambda_1\sim C\varepsilon$), while the \emph{global} projective synchronization gap closes exponentially ($1-\kappa\sim2e^{-1/\varepsilon}$). Local smoothing is easy; global ordering is exponentially hard. The certified spectral bound $\rho(D\Rc)<0.45$ (Theorem~\ref{thm:certified}) shows that at fixed, non-vanishing scale the realized linear gap is far wider than the Birkhoff-theoretic worst case: the degeneration is a singular-limit ($\varepsilon\to0$) phenomenon, not a defect at operating scale. This is the spectral content behind the minimum-sampling requirement of the operational frame: a trajectory must accumulate enough budget for the inflection to be reconstructible.
\end{remark}

\begin{remark}[Closing]
The framework is complete in the following sense. Sections~\ref{sec:prelim}--\ref{sec:axioms} fix the class and reduce accumulation, transition and irreversibility to self-reference and dissipation. Section~\ref{sec:renorm} transports the structure across scales. Section~\ref{sec:operational} instantiates it. Sections~\ref{sec:t1}--\ref{sec:t3} prove the three structural theorems. Sections~\ref{sec:operator}--\ref{sec:nogo} construct the order-parameter fixed point and prove that its location is substrate-specific, not universal. Section~\ref{sec:nogo2} extends the elimination to the deployment level: under interventional monitoring with adversive reception no universal static classifier exists, and verdicts are regime-bound. Sections~\ref{sec:ledger}--\ref{sec:spectral} delimit the rate theory. Every numerical inflection is, by Theorem~\ref{thm:nogo}, a measured quantity; what is universal is the structural grammar: the recurrence of one functional form, read off one cumulative observable, across substrates.
\end{remark}

% ============================================================
\begin{thebibliography}{99}

\bibitem{ashby1956} W.~R. Ashby. \emph{An Introduction to Cybernetics}. Chapman \& Hall, 1956.
\bibitem{banach1922} S.~Banach. Sur les op\'erations dans les ensembles abstraits et leur application aux \'equations int\'egrales. \emph{Fund. Math.}, 3:133--181, 1922.
\bibitem{birkhoff1931} G.~D. Birkhoff. Proof of the ergodic theorem. \emph{Proc. Natl. Acad. Sci. USA}, 17(12):656--660, 1931.
\bibitem{birkhoff1957} G.~Birkhoff. Extensions of Jentzsch's theorem. \emph{Trans. Amer. Math. Soc.}, 85:219--227, 1957.
\bibitem{bobkov2003} S.~G. Bobkov. Spectral gap and concentration for log-concave probability measures. In \emph{Geometric Aspects of Functional Analysis}, Lecture Notes in Math. 1807, pp.~37--43. Springer, 2003.
\bibitem{brascamplieb1976} H.~J. Brascamp and E.~H. Lieb. On extensions of the Brunn--Minkowski and Pr\'ekopa--Leindler theorems. \emph{J. Funct. Anal.}, 22:366--389, 1976.
\bibitem{bushell1973} P.~J. Bushell. Hilbert's metric and positive contraction mappings in a Banach space. \emph{Arch. Ration. Mech. Anal.}, 52:330--338, 1973.
\bibitem{coverthomas2006} T.~M. Cover and J.~A. Thomas. \emph{Elements of Information Theory}, 2nd ed. Wiley, 2006.
\bibitem{crooks1999} G.~E. Crooks. Entropy production fluctuation theorem and the nonequilibrium work relation for free energy differences. \emph{Phys. Rev. E}, 60:2721--2726, 1999.
\bibitem{doeblin1937} W.~Doeblin. Sur les propri\'et\'es asymptotiques de mouvement r\'egis par certains types de cha\^ines simples. \emph{Bull. Math. Soc. Roum. Sci.}, 39:57--115, 1937.
\bibitem{evesonnussbaum1995} S.~P. Eveson and R.~D. Nussbaum. An elementary proof of the Birkhoff--Hopf theorem. \emph{Math. Proc. Cambridge Philos. Soc.}, 117:31--55, 1995.
\bibitem{grunbergmodigliani1954} E.~Grunberg and F.~Modigliani. The predictability of social events. \emph{J. Polit. Econ.}, 62(6):465--478, 1954.
\bibitem{johansson2017} F.~Johansson. Arb: efficient arbitrary-precision midpoint-radius interval arithmetic. \emph{IEEE Trans. Computers}, 66(8):1281--1292, 2017.
\bibitem{kadanoff1966} L.~P. Kadanoff. Scaling laws for Ising models near $T_c$. \emph{Physics}, 2:263--272, 1966.
\bibitem{karlin1968} S.~Karlin. \emph{Total Positivity, Vol.~I}. Stanford University Press, 1968.
\bibitem{lawvere1969} F.~W. Lawvere. Diagonal arguments and cartesian closed categories. In \emph{Category Theory, Homology Theory and their Applications II}, Lecture Notes in Mathematics 92, pp.~134--145. Springer, 1969.
\bibitem{maturana1980} H.~R. Maturana and F.~J. Varela. \emph{Autopoiesis and Cognition}. Reidel, 1980.
\bibitem{merton1948} R.~K. Merton. The self-fulfilling prophecy. \emph{Antioch Rev.}, 8(2):193--210, 1948.
\bibitem{oseledets1968} V.~I. Oseledets. A multiplicative ergodic theorem. \emph{Trans. Moscow Math. Soc.}, 19:197--231, 1968.
\bibitem{perdomo2020} J.~C. Perdomo, T.~Zrnic, C.~Mendler-D\"unner, and M.~Hardt. Performative prediction. In \emph{Proc. 37th Int. Conf. Machine Learning (ICML)}, PMLR 119, pp.~7599--7609, 2020.
\bibitem{prigogine1977} I.~Prigogine. \emph{Self-Organization in Nonequilibrium Systems}. Wiley, 1977.
\bibitem{reimers2019} N.~Reimers and I.~Gurevych. Sentence-BERT: Sentence embeddings using Siamese BERT-networks. In \emph{Proc. EMNLP-IJCNLP}, pp.~3982--3992, 2019.
\bibitem{rump1999} S.~M. Rump. INTLAB --- INTerval LABoratory. In \emph{Developments in Reliable Computing}, pp.~77--104. Kluwer, 1999.
\bibitem{scheffer2009} M.~Scheffer et al. Early-warning signals for critical transitions. \emph{Nature}, 461:53--59, 2009.
\bibitem{seifert2012} U.~Seifert. Stochastic thermodynamics, fluctuation theorems and molecular machines. \emph{Rep. Prog. Phys.}, 75:126001, 2012.
\bibitem{seneta2006} E.~Seneta. \emph{Non-negative Matrices and Markov Chains}, revised ed. Springer, 2006.
\bibitem{simon1954} H.~A. Simon. Bandwagon and underdog effects and the possibility of election predictions. \emph{Public Opin. Q.}, 18(3):245--253, 1954.
\bibitem{wiener1948} N.~Wiener. \emph{Cybernetics}. MIT Press, 1948.
\bibitem{wilson1971} K.~G. Wilson. Renormalization group and critical phenomena. \emph{Phys. Rev. B}, 4:3174--3205, 1971.
\bibitem{yanofsky2003} N.~S. Yanofsky. A universal approach to self-referential paradoxes, incompleteness and fixed points. \emph{Bull. Symb. Log.}, 9(3):362--386, 2003.

\end{thebibliography}

\end{document}
